% --- ОБЩИЕ НАСТРОЙКИ ТЕКСТА ---
\onehalfspacing
\setlength{\parindent}{1.25cm} 

% Настройки шрифтов заголовков (все 14pt, жирные)
\setkomafont{disposition}{\rmfamily\bfseries\normalsize}
\setkomafont{chapter}{\rmfamily\bfseries\normalsize}
\setkomafont{section}{\rmfamily\bfseries\normalsize}
\setkomafont{subsection}{\rmfamily\bfseries\normalsize}

% Настройка оглавления
\setkomafont{chapterentry}{\normalfont\normalsize}
\setkomafont{chapterentrypagenumber}{\normalfont\normalsize}
\KOMAoptions{toc=chapterentrydotfill, numbers=noendperiod}

% --- МАГИЯ ОТСТУПОВ ЗАГОЛОВКОВ (Принудительный сдвиг) ---
% 1. Убираем автоматические выравнивания, которые могут мешать
\renewcommand*{\raggedsection}{} 

% 2. Переопределяем формат печати строк для ВСЕХ разделов (section, subsection и т.д.)
\makeatletter
\renewcommand{\sectionlinesformat}[4]{%
    \hspace{1.25cm}#3#4% Вставляем 1.25см перед номером (#3) и текстом (#4)
}

% 3. Переопределяем формат специально для ГЛАВ (chapter)
\renewcommand{\chapterlinesformat}[3]{%
    \hspace{1.25cm}#2#3% Вставляем 1.25см перед номером (#2) и текстом (#3)
}
\makeatother

% Настройка интервалов (отступы до и после заголовка)
\RedeclareSectionCommand[beforeskip=12pt, afterskip=6pt]{chapter}
\RedeclareSectionCommand[beforeskip=12pt, afterskip=6pt]{section}
\RedeclareSectionCommand[beforeskip=12pt, afterskip=6pt]{subsection}

% --- ТИПОГРАФИКА (Борьба с дырками и переносами) ---
\sloppy 
\hyphenpenalty=1000 
\tolerance=2000

% --- СПИСКИ ---
\setlist{
	leftmargin=0pt,
	nosep,
	wide=\parindent,
	labelsep=0.5em,
	listparindent=\parindent
}
\setlist[itemize,1]{label=--} 
\setlist[enumerate,1]{label=\arabic*.}

% --- ПОДПИСИ ---
\captionsetup{justification=centering, singlelinecheck=false}
\DeclareCaptionLabelFormat{simple}{#1~#2}
\captionsetup{labelformat=simple}
\DeclareCaptionLabelSeparator{emdash}{ --- }
\captionsetup[table]{
	labelsep=emdash,
	justification=raggedright,
	singlelinecheck=false,
	position=top,
	skip=6pt
}

% --- ЛИСТИНГИ ---
\lstset{
    basicstyle=\ttfamily\footnotesize,
    breaklines=true,
    texcl=true,
    frame=none
}

\DefineBibliographyStrings{russian}{
	urlseen = {дата обращения:},
}